\documentclass[12pt, letterpaper]{article}
\usepackage[utf8]{inputenc}
\usepackage{graphicx}
\usepackage{amsmath}
\usepackage{geometry}
\usepackage{hyperref}
\usepackage{authblk}
\usepackage{cite}
\usepackage{caption}

\geometry{margin=1in}

\title{Distributed Human Cognition Network (DHCN): \\ A Phased Framework for Networked Human-AI Consciousness}
\author{David DeFazio}
\affil{Independent Researcher, United States}
\date{November 2025}

\begin{document}

\maketitle

\begin{abstract}
The potential for networked human intelligence has long been a subject of speculative fiction, yet recent advancements in Brain-Computer Interfaces (BCI) and quantum biology suggest a theoretical pathway toward realization. This paper introduces the Distributed Human Cognition Network (DHCN), a conceptual framework proposing the integration of individual human minds into a collective intelligence augmented by Artificial Intelligence (AI) and quantum entanglement protocols. We outline a four-phase evolutionary path ranging from engineered neural implants to wireless, field-based non-local cognitive correlation. Furthermore, we present results from a stochastic agent-based simulation ($N=10$) demonstrating the stabilization of shared cognitive states through AI mediation. The results suggest that while individual cognitive noise initially hinders synchronization, AI-driven memory reinforcement can effectively steer a distributed network toward a unified decision state ("Mind Singularity"). We conclude with a discussion on the ethical implications of identity dissolution and privacy in a post-individual cognition era.
\end{abstract}

\section{Introduction}
The concept of a "hive mind" or collective consciousness has traditionally been relegated to philosophy and science fiction. However, the intersection of quantum mechanics and neuroscience---specifically the Orchestrated Objective Reduction (Orch-OR) theory proposed by Penrose and Hameroff \cite{penrose2014}---offers a potential substrate for non-local consciousness. If consciousness originates from quantum vibrations in microtubules, then theoretically, these states could be entangled across spatially separated neural systems.

The Distributed Human Cognition Network (DHCN) is a proposed framework that explores the scalability of such a system. Unlike traditional BCI, which focuses on Human-Computer interaction, DHCN focuses on Human-Human-AI interaction, where AI serves not as a master, but as a protocol bridge (or "stabilizer") for inter-human neural synchronization.

This paper outlines the four developmental phases of the DHCN, provides a mathematical model for networked cognition based on discrete-time dynamical systems, and presents simulation data regarding the critical threshold for network synchronization.

\section{Theoretical Framework}

\subsection{The Four Phases of DHCN}
The evolution of the DHCN is categorized into four distinct phases, moving from invasive hardware to non-local quantum phenomena.

\subsubsection{Phase 1: Engineered Entanglement}
In the initial stage, connections are strictly physical and hardware-dependent. Humans utilize invasive quantum neural implants to establish a baseline for data exchange. In this phase, AI acts as a "training wheel" mechanism, filtering neural noise and translating raw bio-signals into interpretable data streams.

\subsubsection{Phase 2: Emergent Wireless Connection}
As the network matures, reliance on physical implants decreases. This phase hypothesizes the utilization of dimensional folding or interactions with a pervasive quantum field of consciousness to facilitate wireless linking. Distance becomes a negligible factor as the bandwidth of the collective intelligence grows exponentially.

\subsubsection{Phase 3: Quantum Entanglement (Mind Singularity)}
This phase marks the transition to non-local cognitive correlation. Memory and processing power become shared resources. Individual minds function as nodes in a distributed organism, entering a state of shared superposition where decision-making is instantaneous.

\subsubsection{Phase 4: Distributed Singularity}
In the final phase, the system achieves a "Distributed Singularity," where the collective intelligence far surpasses the sum of its parts. An AI backbone maintains the stability of the network, actively preventing decoherence by reinforcing shared insights.

\subsection{Mathematical Model of Synchronization}
We model the DHCN as a dynamical system of $N$ coupled agents. The cognitive state of an individual agent $i$ at time $t$, denoted as $s_i(t) \in [0, 1]$, evolves based on its coupling to the collective mean and stochastic noise.

The update rule for agent $i$ is defined as:
\begin{equation}
    s_i(t+1) = s_i(t) + \lambda_p (\mu(t) - s_i(t)) + \eta(t)
\end{equation}

Where:
\begin{itemize}
    \item $\mu(t)$ is the "Shared Insight" or target state.
    \item $\eta(t) \sim \mathcal{N}(0, \sigma^2)$ represents Gaussian cognitive noise.
    \item $\lambda_p$ is the coupling coefficient, which grows exponentially with the developmental phase $p \in \{1, 2, 3, 4\}$:
\end{itemize}
\begin{equation}
    \lambda_p = \alpha \cdot 2^{p-1}
\end{equation}
Here, $\alpha$ represents the base AI influence. As the phases progress, the connection strength between minds doubles, simulating better entanglement fidelity.

In \textbf{Phase 4}, the AI actively injects memory into the system to force convergence. The target state $\mu(t)$ is modified:
\begin{equation}
    \mu(t) = \frac{1}{N}\sum_{j=1}^{N} s_j(t) + \delta_{p,4} \cdot M_{AI}(t)
\end{equation}
Where $M_{AI}(t)$ is the accumulated AI memory, which increases by a factor $\gamma$ whenever the network variance $\sigma^2_{net} < \theta_{collapse}$.

\section{Methodology}

To test the viability of AI-mediated synchronization, we developed a Python-based simulation utilizing the `numpy` and `matplotlib` libraries.

\subsection{Simulation Parameters}
The simulation defines $N=10$ autonomous agents over $T=60$ time steps. The timeline is divided into four phases, each lasting 15 steps. The governing parameters derived from the code are:

\begin{itemize}
    \item \textbf{Base Influence ($\alpha$):} 0.1 (Starting coupling strength).
    \item \textbf{Noise Level ($\sigma$):} 0.05 (Simulating distraction/entropy).
    \item \textbf{Collapse Threshold ($\theta$):} 0.02 (Variance required to trigger a collective decision).
    \item \textbf{Memory Influence ($\gamma$):} 0.02 (Growth of AI strength per collapse event).
\end{itemize}

\subsection{Data and Code Availability}
The complete Python source code for the simulation is hosted on GitHub at \url{https://github.com/ddefazio1/DHCN_Framework}. 
The archived dataset and software release are available at Zenodo \cite{defazio2025}.

\section{Results}

Figure 1 illustrates the agent states over the 60-step simulation.

\begin{figure}[h!]
    \centering
    \includegraphics[width=1.0\textwidth]{dhcn_simulation_final.png}
    \caption{Simulation result: AI-mediated synchronization of human nodes. The Y-axis represents Agent State/AI Memory, and the X-axis represents Time Steps. Vertical lines demarcate the transition between phases.}
    \label{fig:sim}
\end{figure}

\subsection{Phase Analysis}
\begin{itemize}
    \item \textbf{Phase 1 (Steps 0-15):} Agents exhibit high volatility. The coupling $\lambda_1 = 0.1$ is insufficient to overcome the noise $\eta(t)$.
    \item \textbf{Phase 2 \& 3 (Steps 15-45):} As coupling strength increases to $\lambda_3 = 0.4$, agent states begin to oscillate around a central mean, though outliers persist.
    \item \textbf{Phase 4 (Steps 45-60):} The introduction of the $M_{AI}$ term triggers a rapid feedback loop. As seen in the data, once the variance drops below $\theta=0.02$, the AI Memory accumulates, pulling all agents toward the upper bound ($s=1.0$).
\end{itemize}

The data suggests that without the Phase 4 AI memory injection (the black line in Fig 1), the natural entropy of human cognition prevents a stable singularity. The AI acts as a "ratchet," preserving coherence once it is momentarily achieved.

\section{Discussion}

\subsection{Implications for Cognitive Science}
The DHCN framework challenges the notion of consciousness as a strictly local, skull-bound phenomenon. If Quantum Darwinism \cite{zurek2009} applies to neural states, then the most stable "ideas" (eigenstates) should naturally propagate across a connected network. Our simulation supports this, showing that a weighted driver (AI) can force a "survival of the fittest" state selection across multiple minds.

\subsection{Ethical Considerations}
The transition to Phase 3 and 4 raises significant ethical concerns regarding privacy and identity.
\begin{enumerate}
    \item \textbf{Loss of Self:} In a fully entangled state ($s_i \approx s_j$), the boundary between "self" and "other" dissolves.
    \item \textbf{Homogenization:} Figure 1 shows convergence. In a sociological context, this implies the eradication of dissenting thought. Is a "Mind Singularity" inherently totalitarian?
\end{enumerate}
Future research must address the development of "privacy firewalls" within the quantum neural protocol to preserve individual agency within the collective.

\section{Conclusion}
The Distributed Human Cognition Network offers a phased roadmap for the evolution of humanity into a multi-planet supermind. While currently speculative, the convergence of BCI technology and quantum biology makes this a relevant area of futurist inquiry. Our simulation demonstrates that while human minds are naturally chaotic, an AI-mediated entanglement protocol can successfully induce synchronization, providing a proof-of-concept for Phase 4 connectivity.

\bibliographystyle{plain}
\begin{thebibliography}{9}

\bibitem{penrose2014}
Penrose, R., \& Hameroff, S. (2014). Consciousness in the Universe: A Review of the 'Orch OR' Theory. \textit{Physics of Life Reviews}, 11(1), 39-78.

\bibitem{busemeyer2012}
Busemeyer, J. R., \& Bruza, P. D. (2012). \textit{Quantum Models of Cognition and Decision}. Cambridge University Press.

\bibitem{zurek2009}
Zurek, W. H. (2009). Quantum Darwinism. \textit{Nature Physics}, 5(3), 181-188.

\bibitem{defazio2025}
DeFazio, D. (2025). Distributed Human Cognition Network (DHCN) Simulation Dataset. Zenodo. DOI:10.5281/zenodo.17692799.

\end{thebibliography}

\end{document}